\documentclass[10pt,final]{beamer}
\mode<presentation>
\usetheme{lucasdeb}
\usepackage{paralist}
\usepackage{wasysym}
\definecolor{darkgreen}{rgb}{0,.7,0}
\definecolor{darkred}{rgb}{.7,0,0}
\newcommand{\Smiley}{{\color{darkgreen}\smiley}}
\newcommand{\Frownie}{{\color{darkred}\frownie}}
\usepackage[french]{babel}
\usepackage[utf8]{inputenc}
\usepackage{helvet}
\usepackage[OT1]{fontenc}
\usepackage{xspace}
\usepackage{verbatim}
\usepackage{comment}
\usepackage{tabularx}
\usepackage{listings}
\lstset{basicstyle=\ttfamily}
\beamertemplatetransparentcovereddynamic
\setbeamercovered{invisible}
\newcommand{\br}{\vskip 1em}
\newcommand{\hbr}{\vskip 0.5em}
\newcommand{\hhbr}{\vskip 0.25em}
\newcommand{\Arr}{\textbf{$\Rightarrow$}\xspace}
\newcommand{\arr}{\textbf{$\rightarrow$}\xspace}
\newcommand{\fixme}[1]{ {\color{red}\footnotesize\bf #1}}
\usepackage{listings}

\renewcommand{\ttdefault}{aett}
\usepackage{aecompl}

\newcommand{\pedp}{\NoAutoSpaceBeforeFDP}

\usebackgroundtemplate{\includegraphics[width=\paperwidth]{figs/dslide8.pdf}}
\logo{\includegraphics[viewport=274 335 360 440,width=0.5cm]{figs/openlogo-nd.pdf}}
\hypersetup{pdftitle={An Introduction to Debian Packaging},bookmarks}
\title{An Introduction to Debian Packaging}
\author[]{Lucas Nussbaum\\{\small\texttt{lucas@debian.org}}}
\institute{\includegraphics[viewport=274 335 360 440,width=1cm]{figs/openlogo-nd.pdf}}
\date{}

\newcommand{\nospace}[1]{{\NoAutoSpaceBeforeFDP{}#1}}%

\begin{document}

\frame{\titlepage}

\section{Introduction}
\subsection{About this document}
\begin{frame}
\frametitle{About this document}

Copyright \copyright 2011 Lucas Nussbaum -- lucas@debian.org
\br

{\small 
This document is free software: you can redistribute it and/or modify
it under either (at your option):
\begin{itemize}
\item the terms of the GNU General Public License as published by
    the Free Software Foundation, either version 3 of the License, or
    (at your option) any later version.\\
    \url{http://www.gnu.org/licenses/gpl.html}
    \hbr
\item the terms of the Creative Commons Attribution-ShareAlike 3.0 Unported License.\\
	\url{http://creativecommons.org/licenses/by-sa/3.0/}
\end{itemize}
}

\br
\fixme{Add pointer to website / git repo}
\end{frame}

\subsection{The Deb package format}
\begin{frame}[fragile]
	\frametitle{Introduction: the Deb package format}
\begin{itemize}
	\item One of the two most common packages format (with RPM)
	\item Used by Debian and its derivatives (including Ubuntu)
	\item A very powerful and convenient way to distribute software to users
		\pause
	\item \texttt{.deb} file: an \texttt{ar} archive
\begin{lstlisting}[basicstyle=\ttfamily\footnotesize]
$ ar tv wget_1.12-2.1_i386.deb
rw-r--r-- 0/0      4 Sep  5 15:43 2010 debian-binary
rw-r--r-- 0/0   2403 Sep  5 15:43 2010 control.tar.gz
rw-r--r-- 0/0 751613 Sep  5 15:43 2010 data.tar.gz
\end{lstlisting}
\begin{itemize}
	\item \texttt{debian-binary}: version of the deb file format, \texttt{"2.0\textbackslash{}n"}
	\item \texttt{control.tar.gz}: metadata about the package\\
		{\small \texttt{\textbf{control}, md5sums, (pre|post)(rm|inst), triggers, shlibs}, \ldots}
	\item \texttt{data.tar.gz}: data files of the package
\end{itemize}
		\pause
\item You could create your \texttt{.deb} files manually\\
	{\footnotesize \url{http://tldp.org/HOWTO/html\_single/Debian-Binary-Package-Building-HOWTO/}}
\item But most people don't do it that way
\end{itemize}
\hbr
\centerline{\textbf{This tutorial: create Debian packages, the Debian way}}
\end{frame}

\begin{frame}
	\frametitle{Ideas dump}
\begin{itemize}
\item Tools you will need
\item General packaging workflow / source package - binary packages

\item Downloading the source
\item Creating a basic Debian source package
\item Debian source package -- files in debian/
\item Example: removing IPv6 support from wget
\item Debhelper7
\item CDBS
\item New source formats
\item Modifying the upstream source Patch systems
\item Per-language policies and packaging helpers
\item Maintainer scripts
\item Building packages -- dpkg-buildpackage pbuilder sbuild
\item Packaging with VCSes debcheckout debcommit debrelease git-buildpackage svn-buildpackage
\item Uploading debsign / dput -- private archive -- reprepro
\item Contributing the package to Debian and Ubuntu

\item Contributing new stuff vs adopting (WNPP, etc) vs co-maintaining
\end{itemize}
\end{frame}

\section{Links}
\begin{frame}{Links}

http://wiki.debian.org/IntroDebianPackaging
http://www.debian.org/devel/
http://www.debian.org/doc/debian-policy/
http://www.debian.org/doc/developers-reference/
http://www.debian.org/doc/maint-guide/
https://wiki.ubuntu.com/PackagingGuide
lp:ubuntu-packaging-guide
debian-mentors
mentors.debian.net
library packaging guide
debhelper 7 slides
\end{frame}

\end{document}
